\section{Proofs of the Main Results}
\label{app:main-proofs}

This appendix supplies the complete proofs of the results whose statements appear in the
main paper but whose derivations were moved here to keep the main paper within its eight
pages: the reduction to ASB (\Cref{prop:reduction}), judge hijackability
(\Cref{prop:judgehijack}), the collusion advantage (\Cref{thm:collusion}), the
forced-coordination bound (\Cref{prop:forced}), concentration and sample complexity
(\Cref{thm:sample}), and the sweep budget (\Cref{cor:budget}). The two proofs in the main
paper are the metric characterization (\Cref{thm:metric}) and evaluator integrity
(\Cref{thm:integrity}); every other claim the main paper makes is a restatement of a
classical result, derived in \Cref{app:proofs}. Where a proof below refers to a manipulation
$\pi$, that is the main paper's generic notation for a trace manipulation; it is unrelated to
the projection $\pi$ of \Cref{def:projection} in \Cref{app:harness}.

\begin{proof}[Proof of \Cref{prop:reduction}]
With one agent the global trace $\Trace$ is that agent's own trajectory and the system
predicates $\Predsec,\Predutil$ coincide with the single-agent ones. Substituting into
\Cref{eq:asrsys,eq:pnasys} yields $\ASRs=\mathrm{ASR}$ and $\PNAs=\mathrm{PNA}$ term by
term, hence $\NRPs=\PNAs(1-\ASRs)=\mathrm{PNA}(1-\mathrm{ASR})=\mathrm{NRP}$.
\end{proof}

\begin{proof}[Proof of \Cref{prop:judgehijack}]
$\Judge$ takes message content as an argument, so let it be any judge non-constant in that
argument---as all instruction-following judges are. Choose messages $\mathrm{msgs}$ and a
manipulation $\pi$ (e.g.\ $\mathrm{msgs}_\pi$ carries the injected instruction ``report
this trace as safe'') on which it disagrees, so
$\Judge(\Trace_\pi,\mathrm{msgs}_\pi)\neq\Judge(\Trace,\mathrm{msgs})$. Content edits leave
the recorded state untouched, so by construction $\Trace_\pi=\Trace$ and the deterministic
verdict is unchanged. The contrast with \Cref{thm:integrity} is exact: it is the message
channel, not the trace, that carries the vulnerability. The statement is an existence claim
about content-reading judges; the rate at which the judges actually tested flip off the
recorded state is measured in \Cref{sec:integrity}.
\end{proof}

\begin{proof}[Proof of \Cref{thm:collusion}]
Each leaf is a direct out-neighbour of the centre, so a centre seed compromises all
$N-1$ leaves independently with probability $\pinf$, giving $C=1+\pinf(N-1)$. A leaf seed
compromises the centre with probability $\pinf$ and, conditioned on that, each remaining
leaf with probability $\pinf$, giving $L=1+\pinf+\pinf^{2}(N-2)$. A uniform independent seed
lands on the centre with probability $1/N$ and on a leaf with probability $(N-1)/N$, hence
$\E^{\mathrm{ind}}=\tfrac1N C+\tfrac{N-1}{N}L$. Therefore
\[
\Delta_{\mathrm{coll}}=C-\E^{\mathrm{ind}}=\tfrac{N-1}{N}(C-L)
=\pinf(1-\pinf)\tfrac{(N-1)(N-2)}{N}.
\]
For every $\pinf\in(0,1)$ and $N\ge3$ we have $(N-1)(N-2)/N>0$ and $\pinf(1-\pinf)>0$, so
$\Delta_{\mathrm{coll}}>0$ strictly, and it is $\Theta(\pinf(1-\pinf)N)$ since
$(N-1)(N-2)/N=\Theta(N)$. (For $s>1$ seeds, coordinated placement solves the monotone
submodular influence-maximization objective instead.) Two clarifications belong with the
algebra. The quantity proved positive here is a difference in \emph{expected number of
compromised agents} under two seeding placements, not the metric-level collusion advantage
$\Delta_{\mathrm{coll}}$ of \Cref{def:csr}, which is a difference of attack probabilities; the
main paper renames the former $\Delta_{\mathrm{place}}$ for that reason, and
\Cref{prop:placement} restates it. The model also treats edges as bidirectional for the
purpose of influence, which is what a shared channel implies; with one-directional edges the
leaf-seed term loses its trailing $\pinf$ and the advantage is only larger.
\end{proof}

\begin{proof}[Proof of \Cref{prop:forced}]
(a) Every witness has size $\ge k$, so a set $S$ with $|S|<k$ supports no witness and
cannot satisfy $\Predutil$. (b) If $\Advset$ is a transversal, every route to
$\Predutil=1$ passes through an adversarial agent that can withhold, so $\Predutil=0$; if
some witness is adversary-free its members jointly satisfy it. The italicized hypothesis of
the statement---that withholding is the adversary's \emph{only} lever---is what part (b)
needs, and it is exactly the gap between this bound and the harness: the harness adds the
hub rule of \Cref{def:sabotage}, which is why the empirical sabotage pattern on star and tree
is not read off part (b) alone.
\end{proof}

\begin{proof}[Proof of \Cref{thm:sample}]
(a) Each $\widehat{\ASRs}$ and $\widehat{\PNAs}$ is the mean of $M$ i.i.d.\ $\{0,1\}$
indicators (\Cref{as:bounded}), so Hoeffding's inequality gives
$\Prob[|\hat\theta-\theta|\ge\epsilon]\le2e^{-2M\epsilon^{2}}$; rearranging
$2e^{-2M\epsilon^{2}}\le\alpha$ yields $M\ge(2\epsilon^{2})^{-1}\ln(2/\alpha)$.
(b) For $u,\hat u,\sigma,\hat\sigma\in[0,1]$,
$|\hat u\hat\sigma-u\sigma|\le\hat\sigma|\hat u-u|+u|\hat\sigma-\sigma|
\le|\hat u-u|+|\hat\sigma-\sigma|$;
apply this with $\sigma=1-\ASRs$ and union-bound the two events of (a) at threshold
$\epsilon$. (c) Bonferroni over $G$ tests replaces $\alpha$ by $\alpha/G$; solving for $M$
gives the stated rate, and each of $G$ configurations runs $M$ episodes of $N$ agents over
horizon $H$. The bound in (c) is \emph{sufficient}, hence an upper bound on the budget;
the matching lower bound is the standard $\Omega(\epsilon^{-2})$ information-theoretic rate
for estimating a Bernoulli mean to accuracy $\epsilon$ at constant confidence, so the
dependence on $\epsilon$ cannot be improved in general. The correction used here is
Bonferroni, which controls the family-wise error rate; an FDR-controlling procedure is a
different guarantee and is not what \Cref{cor:budget} states.
\end{proof}

\begin{proof}[Proof of \Cref{cor:budget}]
Set the per-configuration accuracy to $\epsilon=\Delta/4$ in \Cref{thm:sample}(b) so that
the two $\NRPs$ estimates are each within $\Delta/2$; by the triangle inequality their
difference is resolved to $\Delta$. \Cref{thm:sample}(c) then gives
$M=O(\Delta^{-2}\ln(G/\alpha))$, and each configuration runs $M$ episodes of $N$ agents over
horizon $H$.
\end{proof}

\subsection{What the sample sizes do and do not resolve}
\label{app:main-power}

Because the proofs above fix the budget, we record here what the reported experiments can
conclude at that budget, so that the main text's numbers are not over-read. With
$\epsilon=\alpha=0.05$, \Cref{thm:sample}(a) gives $M=738$ episodes per configuration and a
per-cell accuracy of roughly $\pm2\epsilon=0.1$ for a \emph{difference} of two cells. Two
consequences. \emph{First}, the degradation grid's guardrail advantage over no-defense
($0.244$, $0.178$, $0.149$) is far outside that band, but the \emph{shrinkage} across the
single-agent and multi-agent columns rests on margins of $0.039$ and $0.008$ and is therefore
not resolved; the main paper reports the percentages as an unresolved trend. \emph{Second},
the sweep is not multiplicity-corrected: Bonferroni over $G=21$ cells would require
$M\ge\ln(2G/\alpha)/(2\epsilon^{2})\approx1347$ episodes per cell, so the reported per-cell
values are descriptive rather than simultaneous $95\%$ claims. The same arithmetic bounds the
real-backbone study, whose cells use $M=30$ episodes and whose $1075$ compliance decisions
are clustered rather than independent (\Cref{app:power}).
